\documentclass[10pt,a4paper]{amsart}
\usepackage[margin=19mm]{geometry}
\usepackage{amsmath,amssymb,mathtools}
\usepackage[scr=rsfs,cal=euler]{mathalfa}
\usepackage{xfrac}
\usepackage[unicode,hidelinks,bookmarks=false]{hyperref}
\newcommand{\ie}{i.e.\ }
\newcommand{\cf}{cf.\ }
\newcommand{\resp}{respectively\ }
\newcommand{\Subsection}{\subsection}
\providecommand{\nobreakdash}{\nobreak\hbox{-}}
\setlength{\emergencystretch}{2em}
\multlinegap0pt
\usepackage[shortlabels]{enumitem}
\usepackage{amssymb}
\usepackage{mathtools}
\usepackage{xcolor}
\usepackage{tikz}
\newtheorem{lemma}{Lemma}
\newcommand{\Ee}{\mathbb E}
\newcommand{\Pp}{\mathbb P}
\newcommand{\Z}{\mathbb Z}
\title{Aperiodic random walks: uniform estimates, local limits and invariance principles}
\author{Jiaming Xu}
\date{Source: arXiv:2609.02257v1 (CC BY 4.0)}
\begin{document}
\maketitle

\noindent\emph{Source and licence.} Aperiodic random walks: uniform estimates, local limits and invariance principles. Version arXiv:2609.02257v1 (CC BY 4.0), distributed by arXiv under the Creative Commons Attribution 4.0 International licence, http://creativecommons.org/licenses/by/4.0/, as recorded in the arXiv licence metadata for this exact version and shown on the abstract page https://arxiv.org/abs/2609.02257v1. Full licence notice: \texttt{licenses/arxiv-2609.02257-CC-BY-4.0.txt}.\par\smallskip

\noindent\textit{Editorial scope.} Counted unit: the article's lemma lem:unconditioned (source0.tex lines 149--173). The article's complete original proof of it is delivered with it (source0.tex lines 175--268, one \texttt{proof} environment of 94 lines). No mathematics is written, repaired or re-derived: the statement and the proof are the article's own lines, in the article's own notation, and no retained line is substituted or re-flowed.  Retained uncounted as the source-local context the proof needs: lines 48--90.  Nothing was omitted from any retained span, so the package carries no omission record.  No retained line contains a citation, so the wrapper carries no bibliography and no bibliography entry.  No retained line contains a figure environment, an image-inclusion command or drawing code, so no image is delivered and the package declares no asset dependency.  Editorial formatting only, in addition: an AMS \texttt{amsart} preamble that restates the article's own declarations and macros used by the retained lines, namely the environments \texttt{lemma} on separate counters, together with the standard packages those lines need.  The retained environments print in this document as Lemma 1 in order of appearance, whereas the article prints the same items with the numbering of its own full document.  The complete native source of the article as distributed in the arXiv e-print for this exact version is bundled unchanged as \texttt{sources/209185/source0.tex}.
We consider the following triangular array of integer-valued random
walks. For \(N=1,2,\ldots\), let \(X_{N,1},X_{N,2},\ldots\) be i.i.d. with
law \(\nu_N\) on \(\Z\), and write
\[
 S_N(k)=\sum_{j=1}^k X_{N,j},\qquad
 \sigma_N^2=\Ee X_{N,1}^2 .
\]
We impose the following hypotheses.
\begin{enumerate}[label=\textup{(A\arabic*)},ref=\textup{(A\arabic*)}]
\item \label{ass:centered}
      \(\Ee X_{N,1}=0\) for every \(N\), and
      \(\sigma_N^2\to\sigma^2\in(0,\infty)\).
\item \label{ass:exponential}
      There is \(\eta_0>0\) such that
      \[
        \sup_N \Ee e^{\eta_0|X_{N,1}|}<\infty .
      \]
\item \label{ass:weak}
      \(\nu_N\Rightarrow\nu\), where \(\nu\) is an integer-valued
      probability law.
\item \label{ass:aperiodic}
      The array is uniformly strongly aperiodic: for every
      \(\delta\in(0,\pi)\), there is \(\rho_\delta<1\) such that,
      for all sufficiently large \(N\),
      \[
        \sup_{\delta\le |t|\le\pi}
        \left|\sum_{z\in\Z}e^{itz}\nu_N(z)\right|
        \le \rho_\delta .
      \]
\end{enumerate}
Assumptions~\ref{ass:exponential} and~\ref{ass:weak} imply that, for every
\(0<\eta<\eta_0\),
\begin{equation}
  \sum_{z\in\Z} e^{\eta|z|}
  |\nu_N(z)-\nu(z)|\longrightarrow0 .
  \label{eq:weighted-tv}
\end{equation}
Indeed, weak convergence gives convergence on every finite subset of
\(\Z\), while the two tails are uniformly negligible by
Assumption~\ref{ass:exponential}. Consequently, all moments converge.
Together with Assumptions~\ref{ass:centered} and~\ref{ass:aperiodic}, this
also shows that the limit law is centered, has variance \(\sigma^2\), and is
aperiodic.

\begin{lemma}[Uniform local and local-maximum bounds]
\label{lem:unconditioned}
There are constants \(C,c>0\), independent of \(N,L,z,R\), such that
for all sufficiently large \(N\), \(L,R\in\Z_{\ge0}\), and \(z\in\Z\),
\begin{align}
 \Pp(S_N(L)=z)
 &\le \frac{C}{\sqrt{L+1}}
 \exp\left\{-c\frac{z^2}{L+|z|+1}\right\},       \label{eq:local-bernstein}\\
 \Pp\left(S_N(L)=z,
       \max_{0\le j\le L}|S_N(j)|\ge R\right)
 &\le \frac{C}{\sqrt{L+1}}
 \exp\left\{-c\frac{R^2}{L+R+1}\right\}.         \label{eq:local-max}
\end{align}
Moreover, for some sufficiently large \(N_0\), the uniform lattice
local central limit theorem holds:
\begin{equation}
 \lim_{L\to\infty}\sup_{N\ge N_0}\sup_{z\in\Z}
 \left|
 \sigma_N\sqrt L\,\Pp(S_N(L)=z)
 -\frac1{\sqrt{2\pi}}
  \exp\left\{-\frac{z^2}{2\sigma_N^2L}\right\}
 \right|=0.
 \label{eq:uniform-local-clt}
\end{equation}
\end{lemma}

\begin{proof}
Write
\[
 \phi_N(u):=\Ee e^{iuX_{N,1}},
 \qquad
 \psi_N(t):=\log\Ee e^{tX_{N,1}} .
\]
We first record the uniform local concentration estimate. By
Assumptions~\ref{ass:centered} and~\ref{ass:exponential},
\(\sigma_N^2\) is bounded away from zero and
\(\Ee |X_{N,1}|^3\) is bounded uniformly in \(N\). Hence, for some
\(\delta,c>0\) and all sufficiently large \(N\),
\begin{equation}\label{eq_1}
 |\phi_N(u)|\le e^{-cu^2},\qquad |u|\le\delta .
\end{equation}
On \(\delta\le |u|\le\pi\), Assumption~\ref{ass:aperiodic} gives
\begin{equation}\label{eq_2}|\phi_N(u)|\le\rho<1.\end{equation}
Fourier inversion therefore gives
\[
 \sup_{z\in\Z}\Pp(S_N(L)=z)
 \le
 \frac1{2\pi}\int_{-\pi}^{\pi}|\phi_N(u)|^L\,du
 \le \frac{C}{\sqrt{L+1}} .
\]
The same Fourier decomposition proves
\eqref{eq:uniform-local-clt}.  Indeed, split the Fourier integral into \(|u|\le\delta\) and
\(\delta\le |u|\le\pi\).  In the first region, the change of variables
\(v=\sigma_N\sqrt L\,u\), the uniform Taylor expansion of \(\phi_N\),
and \eqref{eq_1} give \(L^1(dv)\)-convergence to \(e^{-v^2/2}\).
The second region contributes at most \(C\sqrt L\,\rho^L\) by
\eqref{eq_2}.  Fourier inversion then proves
\eqref{eq:uniform-local-clt}, uniformly in \(z\).

We now add the exponential factor by a tilting argument. For
\(|t|\le t_0\), with \(t_0>0\) small, set
\[
 \nu_{N,t}(x):=e^{tx-\psi_N(t)}\nu_N(x),
 \qquad
 \phi_{N,t}(u):=\sum_{x\in\Z}e^{iux}\nu_{N,t}(x).
\]
The tilted laws are not centered. But it suffices to note that under our assumptions, for $|t|\le t_0$, their exponential moments are bounded in a
common neighborhood of the origin, their variances are bounded above
and away from zero, and their characteristic functions have a uniform
gap away from the origin. In particular,
for sufficiently small \(t_0\), there are constants \(c>0\) and
\(\rho<1\) for which \eqref{eq_1}--\eqref{eq_2} hold uniformly in
\(N\) and \(t\). Thus the preceding Fourier bound
applied to \(\nu_{N,t}\) gives
\[
 \sup_{w\in\Z}\Pp_{N,t}(S_N(L)=w)\le \frac{C}{\sqrt{L+1}},
\]
uniformly in \(N,L,t\). Since the original laws are centered and have a
uniform exponential moment,
\[
 \psi_N(t)\le Ct^2,\qquad |t|\le t_0 .
\]
Moreover
\[
 \Pp(S_N(L)=z)
 =
 e^{-tz+L\psi_N(t)}\Pp_{N,t}(S_N(L)=z).
\]
Choose
\[
 t=\operatorname{sgn}(z)\,\alpha\min\{|z|/L,1\}
\]
when \(L\ge1\), with \(\alpha>0\) small enough. If \(|z|\le L\), then
\(-tz+L\psi_N(t)\le -c z^2/L\); if \(|z|>L\), then
\(-tz+L\psi_N(t)\le -c|z|\). Enlarging \(C\) to cover \(L=0\) gives
\eqref{eq:local-bernstein}.

For the maximum estimate without the endpoint, the process
\(\exp\{tS_N(j)-j\psi_N(t)\}\) is a martingale. Stopping it at the first
entrance into \([R,\infty)\) gives
\[
 \Pp\left(\max_{0\le j\le L}S_N(j)\ge R\right)
 \le \exp\{-tR+L\psi_N(t)\}.
\]
With \(t=\alpha\min\{R/L,1\}\), and then applying the same argument to
\(-S_N\), we obtain
\[
 \Pp\left(\max_{0\le j\le L}|S_N(j)|\ge R\right)
 \le C\exp\left\{-c\frac{R^2}{L+R+1}\right\}.
\]

It remains only to keep the endpoint. Let \(m=\lfloor L/2\rfloor\). If
\(|z|\ge R/2\), the bound follows from \eqref{eq:local-bernstein}. If
\(|z|<R/2\), an exit from \([-R,R]\) occurs either during the first half,
or, after reversing the second half from its endpoint \(z\), the reversed
walk has absolute displacement at least \(R/2\). In either case we apply
the preceding maximal bound to the half that contains this displacement
and the uniform \(C(L+1)^{-1/2}\) local concentration bound to the other
half. This proves \eqref{eq:local-max}.
\end{proof}

\end{document}
